% Section 16: limitations and roadmap.
\section{Limitations and roadmap}\label{sec:roadmap}

\subsection{Limitations}
\begin{itemize}
  \item \textbf{Speed.} The dual simplex is about \gmrSpx{} times slower than HiGHS on Netlib
        (geometric mean of per-instance ratios; \ratioSpx{} in shifted geometric mean), and
        branch-and-bound is \ratioMip{} times slower on the MIPLIB instances it solves. The iteration
        drivers are Python; only the kernels are compiled.
  \item \textbf{MILP strength.} \nMip{} of \nMipTot{} MIPLIB~3 instances are proved optimal in
        \tlMip\,s against \nMipRefOpt{} for HiGHS. The MILP case studies do not yet solve. MIP presolve
        only tightens bounds, Gomory mixed-integer is the only cut family, the heuristics are simple, and
        there is no parallel tree search.
  \item \textbf{Presolve.} Only reductions with a simple exact postsolve exist so far
        (\cref{sec:presolve}). Doubleton equations, forcing rows, dominated columns, parallel rows and
        columns, and coefficient tightening are missing, and QPs are not presolved at all.
  \item \textbf{Interior points.} Infeasibility is detected heuristically, not certified. There is no
        crossover to a basic solution, and the QP termination test can accept a point whose
        objective is wrong (\cref{sec:results}).
  \item \textbf{Linear algebra.} FTRAN and BTRAN use dense work vectors, updates are product-form, and
        the Cholesky is up-looking rather than supernodal, so it uses no dense BLAS kernels.
  \item \textbf{Scope.} No MIQP, no second-order cones, no SOS constraints, no callbacks or modelling
        layer beyond the Python API and MPS files.
\end{itemize}

\subsection{Roadmap, in order of expected payoff}
\begin{enumerate}
  \item \textbf{Move the simplex driver into compiled code.} Moving the whole iteration into one Numba
        (or C) loop removes the per-call overhead that dominates small problems. We expect this to be
        the largest single speed-up.
  \item \textbf{Complete presolve}, following the census order: doubleton equations and the
        aggregator, forcing rows, dominated columns, parallel rows and columns, then MIP coefficient
        tightening and probing.
  \item \textbf{MIP:} MIR and knapsack-cover cuts with a cut pool and ageing, feasibility jump and a
        feasibility pump for first incumbents, RINS/RENS-style neighbourhood search, reduced-cost fixing
        and a root restart.
  \item \textbf{Linear algebra:} hyper-sparse FTRAN/BTRAN, a Forrest--Tomlin update, and a supernodal
        Cholesky/$LDL\T$ with dense kernels, which would also make a GPU factorisation possible.
  \item \textbf{Interior points:} a homogeneous self-dual formulation for certified infeasibility,
        crossover to a vertex, and a stricter QP gap test.
  \item \textbf{Concurrent LP:} race dual simplex, interior point and GPU PDLP and take the first answer,
        as cuOpt does. The three engines already exist.
  \item \textbf{Refinery-scale validation:} run the case-study MILPs to proven optimality, and add
        MRPL's own planning models when they become available.
\end{enumerate}

\subsection{Conclusion}
\nirnay{} shows that a solver meeting the independence requirement of SIH26119 can be built: every
factorisation and algorithm is our own, the answers are verified against independent solvers, and on
the standard LP and QP test sets it is correct. It is not yet fast, and its MILP code is far behind the
state of the art. The report documents both the engine and its numerical failures, so the work can be
continued.
